
\subsection*{Abstract}
We introduce the \textbf{Bridge Twin Engine}, a mechanism for autonomous topological evolution in Knowledge Graphs (KGs) based on reasoning experience. While traditional KGs maintain a static structure, the Bridge Twin Engine implements a graph-theoretic analog to biological Long-Term Potentiation (LTP) \cite{hebb1949, bipoo1998}. By materializing "Twin Nodes" to act as structural relays across frequently traversed community boundaries, the system short-circuits latent paths and optimizes its own geometry for future inference. We formalize the potentiation rules based on usage frequency and semantic salience, and provide a pruning mechanism (LTD) to maintain topological sanity. The v2.71.0 release incorporates an atomic re-mapping protocol to synchronize relays with background community re-partitioning, effectively solving the "Zombie Bridge" problem in dynamic environments. As of v2.71.0, the GraphBridgeEngine (Phase 30) performs proactive cross-component bridge synthesis before queries encounter dead ends, scaling to 8,300 active bridges in WebQSP OPT configuration and contributing to H@10=20.84\% (versus 16.59\% without bridges); async bridge synthesis via TaskQueue further decouples bridge updates from active beam traversal.

\subsection*{1. Introduction}
The efficiency of multi-hop reasoning in graphs is fundamentally constrained by the diameter and sparsity of the network. In large-scale KGs, semantically related concepts often reside in disparate community partitions, necessitating high-latency "Bridge" traversals. We propose that a graph should not be a static artifact, but a plastic entity that reshapes its topology to favor successful reasoning chains.

\subsection*{2. Methodology}

\subsubsection*{2.1 The Potentiation Rule}
A structural relay node $v'_{twin}$ is materialized for node $v$ in destination community $C_{dest}$ if:
\begin{enumerate}
  \item  \textbf{Frequency}: Traversal $(u \to v)$ occurs $\geq n_{min}$ times.
  \item  \textbf{Alignment}: $\cos(\vec{e}_v, \vec{c}_{dest}) \geq \sigma$, where $\vec{c}_{dest}$ is the community centroid.
\end{enumerate}


\subsubsection*{2.2 Materialization and Reflection}
Upon potentiation, the engine performs a "Reflection" operation:
\begin{itemize}
  \item Assign $v'_{twin}$ to $C_{dest}$.
  \item Establish a \texttt{BRIDGE\_TWIN} binding edge $(v, v'_{twin})$.
  \item Mirror all of $v$'s edges that terminate in $C_{dest}$ onto $v'_{twin}$.
\end{itemize}


\subsection*{3. Structural Maintenance: The LTD Analog}
To prevent the "Informational Sprawl" of materialized nodes, we implement a temporal decay function:
$$c_t = c_0 \cdot \lambda^{\Delta t}$$
where $\lambda$ is the decay constant and $\Delta t$ is the time since last usage. Edges falling below a confidence threshold $c_{min}$ are pruned during the \textbf{REM Cycle} [Buchorn, 2026].

\subsection*{4. Conclusion}
The Bridge Twin Engine provides a biologically-inspired framework for self-optimizing Knowledge Graphs. By transforming reasoning history into physical graph structure, it enables sub-millisecond multi-hop reasoning on increasingly complex network topologies. In CEREBRUM v2.71.0, proactive bridge synthesis scales to 8,300 active bridges in the WebQSP OPT configuration, contributing directly to a H@10 improvement from 16.59\% to 20.84\% --- a concrete, quantified demonstration of experience-dependent structural plasticity improving reasoning recall.

\medskip\noindent\rule{\linewidth}{0.4pt}\medskip

\section*{5. Recent Advances (v2.51.1 -> v2.71.0)}

The Bridge Twin Engine has been substantially extended since v2.51.1. The following describes the most significant developments.

\textbf{GraphBridgeEngine: Proactive Cross-Component Synthesis (Phase 30).} The original Bridge Twin Engine was reactive --- bridges materialized only after a traversal encountered a community boundary. The \texttt{GraphBridgeEngine} (Phase 30) introduces proactive synthesis: at graph load time and after each community re-partitioning, the engine identifies disconnected or weakly connected components and pre-materializes bridge nodes across their boundaries. This eliminates the cold-start penalty where early queries must "earn" bridges before benefiting from them.

\textbf{Scale: 8,300 Active Bridges on WebQSP.} In the WebQSP OPT configuration, the system maintains approximately 8,300 active bridge records. The impact on recall is measurable: H@10 rises from 16.59\% (without bridge synthesis) to 20.84\% (with GraphBridgeEngine), a relative gain of +25.6\%.

\textbf{Async Bridge Synthesis via TaskQueue (Phase 39).} Prior to Phase 39, bridge materialization occurred synchronously in the beam traversal hot path, adding latency to queries that triggered new potentiation events. Phase 39 decouples this via a \texttt{TaskQueue}: potentiation events are enqueued and processed by a background worker. Active traversals are unaffected by bridge write operations, and the community map swap post-rebalance remains atomic.

\textbf{Post-Rebalance Bridge Pruning (Phase 19).} After the \texttt{GlobalRebalancer} triggers a background DSCF/TSC re-run and performs the atomic community map swap, a post-rebalance hook notifies \texttt{BridgeTwinEngine} to prune stale bridge records whose source or destination community assignments have changed. This prevents "Zombie Bridges" accumulating after topology shifts.

\textbf{Benchmark Impact.}

\begin{table}[h]
\small\centering
\begin{tabular}{ll}
\toprule
Configuration & H@10 (WebQSP) \\
\midrule
No bridge synthesis & 16.59% \\
With GraphBridgeEngine (OPT) & 20.84% \\
Gain & +25.6% relative \\
\bottomrule
\end{tabular}
\end{table}
